\documentclass[11pt,a4paper]{article}
\usepackage[margin=1in]{geometry}
\usepackage[T1]{fontenc}
\usepackage[utf8]{inputenc}
\usepackage{lmodern}
\usepackage{amsmath,amssymb,booktabs,microtype}
\usepackage[numbers,sort&compress]{natbib}
\usepackage{xurl}
\usepackage[colorlinks=true,linkcolor=blue,citecolor=blue,urlcolor=blue]{hyperref}
\setlength{\emergencystretch}{2em}
\newcommand{\E}{\mathbb{E}}
\newcommand{\R}{\mathbb{R}}
\newcommand{\placeholder}[1]{\par\medskip\noindent\textit{[Placeholder: #1]}\par\medskip}
\title{The Role of Dynamical Structure in Cross-Modal Representation Learning}
\author{Alex Martinelli}
\date{Working manuscript --- November 2026}

\begin{document}
\maketitle

\begin{abstract}
Independently trained models can develop similar representation geometries, but shared information alone does not determine the geometry selected by learning. This study asks whether temporal locality relative to fixed observation geometry improves finite-sample cross-modal alignment under a specified self-supervised objective. We propose a controlled comparison between a local Markov process and a permutation-conjugated process with unchanged observation channels. The construction preserves the stationary cross-modal observation distribution, latent transition entropy, and complete transition spectrum while changing which transitions are local in observation space. Small encoders are trained separately on independently sampled unimodal temporal pairs, without cross-modal losses or latent-state targets. Evaluation distinguishes geometric alignment, collapse, recovery of accessible factors, and approximation of transition-specific spectral coordinates. Four experiments cover construction audits, data efficiency, measurement noise and aliasing, and low-resolution image--sensor observations. The intended contribution is a bounded test of the interaction between dynamics, observations, and learning bias, not a universal physical explanation of representational convergence. This manuscript specifies the protocol; empirical results and conclusions remain to be supplied.
\end{abstract}

\noindent\textbf{Keywords:} representation alignment; temporal locality; self-supervised learning; spectral methods; multimodal learning.

\section{Introduction}
The Platonic Representation Hypothesis (PRH) proposes that increasingly capable models may converge toward a shared statistical representation of reality \citep{huh2024}. This motivates a distinction between the information available to different models and the geometric organization of their learned features. Two encoders can retain equivalent information through invertible transformations without preserving distances, inner products, or an orthogonal correspondence. Thus, observing a common latent process is not sufficient to guarantee geometric alignment.

Physical systems motivate the study of temporal regularity because nearby states often evolve through structured interactions. However, an experiment comparing ``physical'' data with randomized data can simultaneously alter information availability, uncertainty, temporal dependence, and approximation difficulty. Such an experiment cannot attribute its outcome to physics as a separate cause. Conversely, imposing a common latent-state prediction target introduces supervision and cannot establish convergence from observations alone.

We replace that broad question with a specific one: does the compatibility between temporal neighborhoods and fixed sensor geometry influence alignment at finite data and optimization budgets? Our intervention preserves several potential alternative explanations exactly. It does not preserve every temporal statistic, because those statistics carry the mechanism being investigated. This distinction between controlled quantities and intentionally altered quantities is central to the design.

The proposed contributions are: (i) an auditable isospectral comparison with identical instantaneous cross-modal information; (ii) an evaluation separating geometric agreement from accessible-factor recovery and collapse; and (iii) a low-compute implementation suitable for independently reproducible experiments. The distributional identities are elementary controls, not claimed new theorems. Any empirical novelty depends on the eventual findings and a continued review of related work.

\section{Related Work}
\paragraph{Convergence and identifiability.}
PRH discusses shared reality, task generality, capacity, and simplicity bias \citep{huh2024}. Balestriero and LeCun establish an idealized JEPA convergence result under assumptions including injective renderings, matched positive-pair structure, and minimizers unique up to rotation \citep{balestriero2026}. Such sufficient conditions are not a statement that arbitrary objectives or lossy sensors must align. Linear identifiability \citep{roeder2021}, contrastive recovery of generative factors \citep{zimmermann2021}, content--style identification \citep{vonkugelgen2021}, and multimodal contrastive identifiability \citep{daunhawer2023} further clarify the importance of assumptions and equivalence classes. Joint multimodal objectives should not be confused with independently trained encoders. General unsupervised disentanglement is not identifiable without additional inductive assumptions \citep{locatello2019}.

\paragraph{Temporal and spectral representation learning.}
Slow feature analysis learns slowly varying features subject to nontriviality constraints \citep{wiskott2002}; contrastive predictive coding uses temporal prediction to learn representations \citep{oord2018}. Laplacian eigenmaps and diffusion maps connect data relationships to spectral coordinates \citep{belkin2003,coifman2006}. Modern self-supervised objectives admit related spectral interpretations \citep{haochen2021,balestriero2022}. Eigenspace perturbation theory explains why separation at a selected spectral boundary matters \citep{yu2015}. These established connections motivate, rather than constitute, the contribution proposed here.

\paragraph{Objectives and measurement.}
SimCLR, Barlow Twins, and VICReg provide influential approaches to learning from positive pairs and controlling degeneracy \citep{chen2020,zbontar2021,bardes2022}. We use a small temporal adaptation of VICReg, implemented in PyTorch \citep{paszke2019} with AdamW \citep{loshchilov2019}. Alignment is measured using linear centered kernel alignment (CKA), representational similarity analysis, and orthogonal Procrustes \citep{kornblith2019,kriegeskorte2008,schonemann1966}. These measurements capture different aspects of geometry; none alone establishes recovery of physical laws.

\section{Research Question and Hypotheses}
\subsection{Primary question}
\begin{quote}
Does temporal locality relative to fixed observation geometry improve finite-sample alignment between independently trained temporal encoders, compared with isospectral scrambled dynamics, when the stationary observation distribution and latent transition entropy are held fixed?
\end{quote}
The term ``structure'' here refers to temporal locality relative to specified sensors, not physicality in general. The local and scrambled processes are isomorphic as abstract Markov chains. They differ in their relationship to the observation maps.

\subsection{Prespecified comparisons}
For training budget $N$, let $A_{r,N,c}$ be held-out CKA in independent world $r$ under condition $c$, averaged over initialization/data repeats within that world. The primary estimand is
\begin{equation}
\Delta_N=\E_r[A_{r,N,\mathrm{local}}-A_{r,N,\mathrm{scrambled}}].
\end{equation}
\textbf{H1:} $\Delta_N>0$ at the prespecified finite budgets under the baseline sensor and training configuration. A null result, reversal, or uncertainty spanning practically relevant values does not support this directional hypothesis. A positive result is conditional on the tested learner and observations.

\textbf{H2:} The local--scrambled contrast depends on observation quality and observability. This is a secondary interaction question, not a claim that additional dynamics can recover information absent from a single-frame input.

\textbf{H3:} The direction of a baseline effect persists when one numerical sensor is replaced by low-resolution images and its MLP by a small CNN. This is a synthetic robustness question, not external validity on natural modalities.

\section{Controlled Generative Framework}
\subsection{States and dynamics}
Let $Z_t\in\{0,\ldots,S-1\}$ with $S=32$, and attach circular coordinates
\begin{equation}
c(s)=(\cos(2\pi s/S),\sin(2\pi s/S)).
\end{equation}
Define the local transition matrix, with indices modulo $S$, by
\begin{equation}
P_{ss}=\tfrac12,\qquad P_{s,s-1}=P_{s,s+1}=\tfrac14.
\end{equation}
For a random permutation $\pi$, define $P^{\pi}_{ij}=P_{\pi(i),\pi(j)}$. Both matrices are symmetric and doubly stochastic, so the uniform distribution is stationary. Permutation conjugation preserves eigenvalues, row entropy, and self-transition probabilities. An additional IID control uses $P^0_{ij}=1/S$; it shares the stationary distribution but is not entropy matched.

\subsection{Observation channels}
Each modality observes $X^{(m)}=g_m(Z)+\epsilon_m$, with independent zero-mean Gaussian measurement noise. Numerical sensors apply independently seeded affine projections and a hyperbolic tangent to the circular first and second harmonics, followed by fixed per-channel scaling. These fixed scaling constants belong to the simulator, not an estimated preprocessing fit on test examples. No observation map is changed between local and scrambled conditions within a world.

At stationarity, the instantaneous cross-modal law is
\begin{equation}
p_c(x_a,x_b)=\frac1S\sum_{s=0}^{S-1}p_a(x_a\mid s)p_b(x_b\mid s),
\label{eq:matched}
\end{equation}
which is independent of the transition condition $c$. Instantaneous cross-modal mutual information is therefore equal across conditions by construction, without estimating continuous-variable entropy. This does not imply equality of observation-level temporal mutual information.

For the local and scrambled conditions,
\begin{equation}
H(Z_{t+1}\mid Z_t)=-\frac1S\sum_{i,j}P_{ij}\log P_{ij},\qquad
I(Z_t;Z_{t+1})=\log S-H(Z_{t+1}\mid Z_t)
\end{equation}
are also identical. Both equalities concern latent states and do not imply equal computational learnability from noisy observations.

\subsection{What changes, and what does not}
Temporal locality is summarized by
\begin{equation}
\ell(P)=\frac1S\sum_{i,j}P_{ij}\|c(i)-c(j)\|^2.
\end{equation}
Sensor roughness replaces $c$ with the flattened sensor output and averages over its coordinates. Scrambling generally increases displacement relative to the fixed observation geometry; the realized values are recorded rather than assumed. It changes which observations are temporal neighbors and can change approximation difficulty for a smooth finite network. That is part of the treatment, not a separately eliminated confound.

If the sensor tables are also permuted, $g'_m(i)=g_m(\pi(i))$, the intervention becomes a pure relabeling with an unchanged observation process. This co-permutation is an exact implementation check. It is not the primary experimental condition.

\section{Learning and Evaluation Protocol}
\subsection{Sampling and independence}
Training data consist of independently sampled stationary length-two trajectories: draw $Z\sim\mathrm{Unif}(S)$ and $Z'\sim P(Z,\cdot)$. Modalities receive independent state-pair and noise draws. This removes long-trajectory mixing and overlapping-window leakage from the main comparison. The encoders never receive state IDs, paired cross-modal targets, or one another's embeddings. They share the objective family and underlying transition law, so their independence is procedural rather than absence of common assumptions.

Calibration and test observations are separate balanced draws over the same finite state support, using eight and sixteen observations per state respectively. Consequently evaluation tests fresh measurement noise and independent samples, not unseen-state extrapolation. Calibration pairs are used only after training to fit a restricted alignment map and fixed linear probes.

\subsection{Encoders and objective}
Numerical encoders are two-hidden-layer MLPs with width 64, SiLU nonlinearities, and output dimension $d=4$. The image encoder uses two small convolutional blocks and a width-64 linear layer. Models are initialized independently; architecture and parameter count are fixed across dynamics conditions within each modality. There is no separate projector or shared predictor.

For two minibatch embedding matrices $U,V$, training minimizes
\begin{equation}
\mathcal L=25\mathcal L_{\mathrm{inv}}+25\mathcal L_{\mathrm{var}}+\mathcal L_{\mathrm{cov}},
\end{equation}
where $\mathcal L_{\mathrm{inv}}$ is mean squared difference, $\mathcal L_{\mathrm{var}}$ averages $\max(0,1-\sqrt{\mathrm{Var}(U_j)+10^{-4}})$ over both views and coordinates, and $\mathcal L_{\mathrm{cov}}$ sums squared off-diagonal sample covariances of both views divided by $d$. This is a temporal VICReg-style objective \citep{bardes2022}, not a theorem-equivalent hard whitening constraint.

The reference configuration uses batch size 128, AdamW learning rate $10^{-3}$, weight decay $10^{-4}$, and gradient-norm clipping at 5. The full configured grid uses 800 updates, budgets $N\in\{512,2048,8192\}$, five independent sensor/permutation worlds, and two initialization/data repeats per world. Noise/aliasing comparisons use $N=2048$. Small smoke and pilot configurations are for execution checks and exploration. Fixed updates imply different sample reuse across budgets; this is not an asymptotic sample-complexity analysis. Runtime and parameter counts are recorded, not presumed equivalent across architectures.

\subsection{Spectral reference and its limits}
For an unrestricted state representation $H\in\R^{S\times d}$ with $\mathbf1^TH=0$ and $H^TH/S=I$, the temporal squared-error objective satisfies
\begin{equation}
\E\|h(Z)-h(Z')\|^2=2d-\frac2S\operatorname{tr}(H^TPH).
\end{equation}
Its optimum spans the leading nonconstant eigenvectors of $P$, an established spectral fact \citep{wiskott2002,belkin2003,balestriero2022}. Uniqueness of the selected subspace requires a positive boundary gap $\delta_d=\lambda_d-\lambda_{d+1}$, with eigenvalues ordered from $\lambda_0=1$. The mixing gap $1-\lambda_1$ is different. We choose $d=4$ to include complete repeated-eigenvalue blocks of the cycle. The IID control has no unique preferred nonconstant subspace; spectral recovery scores are undefined there. The oracle uses the full transition matrix and is not a learned model.

\subsection{Measurements}
The primary endpoint is sample-level linear CKA \citep{kornblith2019}. Supporting measurements include an orthogonal Procrustes map fitted on calibration observations only, state-centroid RSA \citep{kriegeskorte2008}, embedding variance, and entropy effective rank of covariance eigenvalues. Procrustes permits centering, one scalar normalization per modality, and an orthogonal map, not arbitrary invertible whitening. State-centroid metrics use latent IDs to average noise and are explicitly oracle-assisted diagnostics.

Fixed ridge probes assess circular coordinates, doubled-angle quotient coordinates, transition-specific spectral coordinates, and the conditional mean of next-state circular coordinates. Probe regularization is fixed at $10^{-3}$. Probe success establishes accessibility to this readout, not causal identification. Original circular coordinates favor the local topology; transition-specific spectral targets are reported separately to expose that distinction. Random-initialization CKA and shuffled evaluation correspondences diagnose pre-existing geometry and measurement artifacts.

\subsection{Noise, aliasing, and inference}
Baseline numerical noise has standard deviation 0.05, with a higher-noise condition at 0.3. An aliased second sensor maps $s$ and $s+S/2$ identically. A single-frame encoder cannot distinguish these states; doubled-angle coordinates specify an accessible quotient. Matching in Equation~\ref{eq:matched} holds within each observation setting, not between different noise or aliasing settings. Image noise is fixed at 0.03.

All runs are retained. Mean embedding variance below $10^{-6}$ or effective rank below $d/2$ triggers a diagnostic collapse flag, not exclusion. Degenerate metrics remain missing rather than being replaced by favorable values. Repeats are averaged within world, followed by paired bootstrap intervals over independent worlds. No interval is reported with fewer than five valid worlds; five-world intervals are exploratory and may be unstable. Observations and initialization repeats are not counted as independent worlds. Multiple budgets and secondary endpoints are reported in full without interpreting exploratory intervals as multiplicity-adjusted tests. A definitive analysis requires a frozen protocol, adequate replication, and disclosure of all pilot-driven changes.

\section{Experiments}
\subsection{Experiment 1: Construction and Measurement Audit}
\placeholder{Insert the control-audit output: matched distributions, entropy and spectra, locality diagnostics, co-permutation checks, and metric sanity tests.}
\subsection{Experiment 2: Finite-Sample Alignment}
\placeholder{Insert local, scrambled, and IID learning comparisons across data budgets, including paired contrasts, initialization baselines, optimization diagnostics, uncertainty, and collapse rates.}
\subsection{Experiment 3: Noise and Partial Observability}
\placeholder{Insert within-setting locality contrasts for noisy and aliased sensors, distinguishing full-coordinate and observable-quotient recovery.}
\subsection{Experiment 4: Low-Resolution Image--Sensor Validation}
\placeholder{Insert CNN--MLP comparisons, measured parameter counts and runtime, geometry metrics, and probe diagnostics.}

\section{Results}
\placeholder{Add the completed empirical results and figures. Do not treat software smoke tests or construction identities as evidence for H1--H3.}

\section{Discussion}
\placeholder{Interpret supported and unsupported hypotheses, alternative explanations, dependence on objective and architecture, and uncertainty. Separate alignment from recovery and generalization.}

\section{Limitations and Scope}
The state space is finite, low dimensional, and synthetic. Isospectral scrambling changes the compatibility of sensor geometry and temporal transitions, not physicality or abstract graph topology. Gaussian noise and exact aliasing represent limited forms of sensor uncertainty. The main encoder processes individual observations rather than histories. Fixed-update comparisons conflate neither conditions nor budgets with extra updates, but may still be optimization limited. The practical loss need not reach its population spectral counterpart, and high initial similarity may leave little room for an alignment advantage.

There is no natural-image, audio, language, or real-world sensor benchmark here, and no claim of out-of-distribution generalization. Matching stationary information does not match temporal information in observations, nor should it eliminate the proposed mechanism. Confirming an effect would not establish a universal law of representation learning; failing to find one would not refute PRH. The intended empirical contribution remains contingent on results and prior-work comparison.

\section{Reproducibility and Data Availability}
The implementation provides four runnable notebooks, NumPy-generated datasets, small PyTorch models, explicit random seeds, tests of construction identities, and checkpointed training. Source, configurations, dependency versions, device information, embeddings, and diagnostics are recorded. No external dataset, private conversation, pretrained model, or proprietary data are needed. Cloud execution starts with an inexpensive smoke profile and can be expanded after runtime measurement. The repository is available at \url{https://github.com/AlexM3020/The-Role-of-Dynamical-Structure-in-Cross-Modal-Representation-Learning}.

\section{Conclusions}
\placeholder{Write conclusions only after completing and interpreting the experiments. State the tested scope, evidence strength, and unresolved questions without claiming proof of universal convergence.}

\begin{thebibliography}{99}
\bibitem{huh2024} M. Huh, B. Cheung, T. Wang, and P. Isola. The Platonic Representation Hypothesis. ICML, 2024. \url{https://arxiv.org/abs/2405.07987}.
\bibitem{balestriero2026} R. Balestriero and Y. LeCun. Proving and Improving the Platonic Representation Hypothesis with JEPAs. Preprint, 2026. \url{https://arxiv.org/abs/2610.39302}.
\bibitem{roeder2021} G. Roeder, L. Metz, and D. Kingma. On Linear Identifiability of Learned Representations. ICML, PMLR 139:9030--9039, 2021. \url{https://proceedings.mlr.press/v139/roeder21a.html}.
\bibitem{zimmermann2021} R. S. Zimmermann, Y. Sharma, S. Schneider, M. Bethge, and W. Brendel. Contrastive Learning Inverts the Data Generating Process. ICML, PMLR 139:12979--12990, 2021. \url{https://proceedings.mlr.press/v139/zimmermann21a.html}.
\bibitem{vonkugelgen2021} J. von K\"ugelgen, Y. Sharma, L. Gresele, W. Brendel, B. Sch\"olkopf, M. Besserve, and F. Locatello. Self-Supervised Learning with Data Augmentations Provably Isolates Content from Style. NeurIPS, 2021. \url{https://arxiv.org/abs/2106.04619}.
\bibitem{daunhawer2023} I. Daunhawer, A. Bizeul, E. Palumbo, A. Marx, and J. E. Vogt. Identifiability Results for Multimodal Contrastive Learning. ICLR, 2023. \url{https://arxiv.org/abs/2303.09166}.
\bibitem{locatello2019} F. Locatello, S. Bauer, M. Lucic, G. R\"atsch, S. Gelly, B. Sch\"olkopf, and O. Bachem. Challenging Common Assumptions in the Unsupervised Learning of Disentangled Representations. ICML, PMLR 97:4114--4124, 2019. \url{https://proceedings.mlr.press/v97/locatello19a.html}.
\bibitem{wiskott2002} L. Wiskott and T. J. Sejnowski. Slow Feature Analysis: Unsupervised Learning of Invariances. Neural Computation, 14:715--770, 2002. \url{https://doi.org/10.1162/089976602317318938}.
\bibitem{oord2018} A. van den Oord, Y. Li, and O. Vinyals. Representation Learning with Contrastive Predictive Coding. Preprint, 2018. \url{https://arxiv.org/abs/1807.03748}.
\bibitem{belkin2003} M. Belkin and P. Niyogi. Laplacian Eigenmaps for Dimensionality Reduction and Data Representation. Neural Computation, 15:1373--1396, 2003. \url{https://doi.org/10.1162/089976603321780317}.
\bibitem{coifman2006} R. R. Coifman and S. Lafon. Diffusion Maps. Applied and Computational Harmonic Analysis, 21:5--30, 2006. \url{https://doi.org/10.1016/j.acha.2006.04.006}.
\bibitem{haochen2021} J. Z. HaoChen, C. Wei, A. Gaidon, and T. Ma. Provable Guarantees for Self-Supervised Deep Learning with Spectral Contrastive Loss. NeurIPS, 2021. \url{https://arxiv.org/abs/2106.04156}.
\bibitem{balestriero2022} R. Balestriero and Y. LeCun. Contrastive and Non-Contrastive Self-Supervised Learning Recover Global and Local Spectral Embedding Methods. NeurIPS, 2022. \url{https://arxiv.org/abs/2205.11508}.
\bibitem{yu2015} Y. Yu, T. Wang, and R. J. Samworth. A Useful Variant of the Davis--Kahan Theorem for Statisticians. Biometrika, 102:315--323, 2015. \url{https://arxiv.org/abs/1405.0680}.
\bibitem{chen2020} T. Chen, S. Kornblith, M. Norouzi, and G. Hinton. A Simple Framework for Contrastive Learning of Visual Representations. ICML, 2020. \url{https://arxiv.org/abs/2002.05709}.
\bibitem{zbontar2021} J. Zbontar, L. Jing, I. Misra, Y. LeCun, and S. Deny. Barlow Twins: Self-Supervised Learning via Redundancy Reduction. ICML, 2021. \url{https://arxiv.org/abs/2103.03230}.
\bibitem{bardes2022} A. Bardes, J. Ponce, and Y. LeCun. VICReg: Variance-Invariance-Covariance Regularization for Self-Supervised Learning. ICLR, 2022. \url{https://arxiv.org/abs/2105.04906}.
\bibitem{kornblith2019} S. Kornblith, M. Norouzi, H. Lee, and G. Hinton. Similarity of Neural Network Representations Revisited. ICML, 2019. \url{https://arxiv.org/abs/1905.00414}.
\bibitem{kriegeskorte2008} N. Kriegeskorte, M. Mur, and P. A. Bandettini. Representational Similarity Analysis---Connecting the Branches of Systems Neuroscience. Frontiers in Systems Neuroscience, 2:4, 2008. \url{https://doi.org/10.3389/neuro.06.004.2008}.
\bibitem{schonemann1966} P. H. Sch\"onemann. A Generalized Solution of the Orthogonal Procrustes Problem. Psychometrika, 31:1--10, 1966. \url{https://doi.org/10.1007/BF02289451}.
\bibitem{paszke2019} A. Paszke et al. PyTorch: An Imperative Style, High-Performance Deep Learning Library. NeurIPS, 2019. \url{https://arxiv.org/abs/1912.01703}.
\bibitem{loshchilov2019} I. Loshchilov and F. Hutter. Decoupled Weight Decay Regularization. ICLR, 2019. \url{https://arxiv.org/abs/1711.05101}.
\end{thebibliography}
\end{document}
